\documentclass{article}

\PassOptionsToPackage{numbers, compress}{natbib}

\usepackage[eandd]{neurips_2026}

\usepackage[utf8]{inputenc}
\usepackage[T1]{fontenc}
\usepackage{hyperref}
\usepackage{url}
\usepackage{booktabs}
\usepackage{amsfonts}
\usepackage{amsmath}
\usepackage{nicefrac}
\usepackage{microtype}
\usepackage{graphicx}
\usepackage{to-be-determined}
\usepackage{xcolor}
\usepackage{placeins}
\newcommand{\todo}[1]{\textcolor{red}{\textbf{TODO:} #1}}

\title{When Alignment Hurts: A Controlled Benchmark for Ranking-Preserving 3D Face Metrics}

% For double-blind submission, the NeurIPS style handles anonymization.
% Replace with real authors only for camera-ready or nonanonymous submission.
\author{%
  Anonymous Authors \\
  Anonymous Institution \\
  \texttt{anonymous@example.com}
}

\begin{document}

\maketitle

\begin{abstract}
The evaluation of 3D face reconstruction is typically based on measuring
geometric distances between reconstructed and ground-truth faces after rigid or
non-rigid alignment. These alignment steps are useful for improving surface overlap, but they introduce a non-uniform bias that distorts the overall similarity structure. As a result, error values become inconsistent across different mesh templates, resolutions, and topologies. We study this problem as a metric meta-evaluation task. Using topology-consistent synthetic faces with known correspondences, we compute true geometric errors and test whether registration-based pipelines preserve the correct ranking of identity differences when topology, sampling, support, and rigid perturbations change. Using a controlled cross-topology benchmark (REMESH), we show that rigid and non-rigid ICP compress the range of distances between subjects. They improve pairwise local surface alignments but reduce identity-related differences, leading to worse ranking consistency. This suggests that reliable 3D face evaluation should separate pairwise alignment quality from identity-aware similarity. To address this, we propose an encoder-only, mesh-agnostic learned metric that maps 3D face meshes into a shared latent space designed to preserve ground-truth dissimilarities using stress, ranking, and identity-consistency objectives. At test time, each mesh is encoded independently and compared in this latent space, without pairwise ICP, non-rigid registration, landmarks, or dense correspondence. On the REMESH benchmark, the proposed metric consistently preserves ground-truth rankings better than standard pipelines.



%The evaluation of 3D face reconstruction is typically based on computing geometric distances between reconstructed and ground-truth faces after rigid or non-rigid alignment. Although such procedures are needed to optimize surface overlap for error computation, we found that they introduce a non-uniform sample bias that severely breaks the global similarity structure, making error values inconsistent across mesh templates, resolutions, and topologies. We study this issue as a metric meta-evaluation problem: given topology-consistent synthetic faces with known exact correspondences from which a real geometric error can be computed, we evaluate the extent to which registration-based error estimation pipelines preserve the ground-truth ranking of identity-level geometric dissimilarities under changes in mesh representation and pose. Using a controlled cross-topology benchmark (REMESH), we show that rigid ICP and non-rigid ICP compress inter-subject distance distributions: they improve local geometric fit while suppressing identity-discriminative morphology, thereby degrading ranking fidelity. These results suggest that robust 3D face evaluation should distinguish pairwise alignment quality from globally comparable identity-aware similarity. To address this issue, we propose an encoder-only, mesh-agnostic learned metric that maps 3D face meshes to a shared latent space designed to preserve correspondence-based ground-truth dissimilarities using stress, ranking, and identity-consistency objectives. At inference time, each mesh is encoded independently and compared in latent space, without pairwise ICP, non-rigid registration, landmarks, or dense correspondence estimation. On the controlled REMESH benchmark, the proposed metric consistently preserves ground-truth rankings better than standard pipelines. 

%%We further show that registration-based pipelines can compress
%%inter-subject distance distributions: they improve local geometric fit while
%%suppressing identity-discriminative morphology, thereby degrading ranking
%%fidelity. These results suggest that robust 3D face evaluation should distinguish
%%pairwise alignment quality from globally comparable identity-aware similarity.
\end{abstract}

\section{Introduction}
Reconstructing 3D faces from single or multiple images is a long-standing problem in computer vision. Recent methods can produce increasingly realistic and detailed facial geometry, even under challenging imaging conditions~\cite{3ddfav2, 3di, deca, deep3d, mica, synergynet, emoca, inorig}. However, as these methods become more diverse in template, topology, resolution, and preprocessing, one key question becomes more critical: \emph{how should we measure reconstruction quality when predicted meshes are not directly comparable?}
%Reconstructing 3D faces from single or multiple images is a long-standing problem in computer vision, with recent methods producing increasingly plausible and detailed facial geometry across challenging imaging conditions~\cite{3ddfav2, 3di, deca, deep3d, mica, synergynet, emoca, inorig}. As reconstruction methods become more diverse in template, topology, resolution, and preprocessing, however, a basic question becomes increasingly important: \emph{how should reconstruction quality be measured when predicted meshes are not directly comparable?}

Standard evaluation protocols, \textit{e.g.}~\cite{now,REALY}, address this question by computing geometric distances, typically the Chamfer distance, after alignment and preprocessing steps. These steps may include rigid alignment, rigid ICP, non-rigid ICP, cropping, landmark initialization, or point-to-surface distance computation. The goal is to make different meshes comparable by enforcing dense correspondence. While these steps are useful, they introduce an important limitation. They optimize pairwise surface alignment, but not the global similarity structure across all meshes; as a result, similarities and differences across the full set of identities may not be preserved. In particular, alignment can reduce geometric differences that carry identity information, such as facial depth, profile shape, or local protrusions. Thus, better pairwise fit does not necessarily mean better identity-aware ranking.
%Standard evaluation protocols \textit{e.g.}~\cite{now,REALY} answer this question by computing a geometric distance (typically the Chamfer distance) after some form of pre-processing and alignment procedures. Rigid alignment, ICP, non-rigid ICP, cropping, landmark initialization, and point-to-surface distances are used to make heterogeneous meshes comparable and serve to establish a dense point-correspondence before measuring error. These procedures are useful, but they also introduce a subtle problem: they optimize pairwise surface fit, not necessarily the global similarity structure required to compare many subjects consistently. In particular, an alignment that reduces the distance between two individual meshes may also suppress identity-discriminative morphology, such as facial depth, profile shape, or local protrusions. Therefore, a better pairwise geometric fit does not necessarily imply a better identity-aware ranking fidelity.

This distinction is important because the true geometric error between two arbitrary 3D face meshes is generally unknown~\cite{10448898}. A well-defined error exists only under strict conditions: meshes must share the same topology, discretization, and canonical alignment. In this setting, distances can be computed directly between corresponding vertices, for example using point-wise RMSE. Synthetic or topology-consistent 3D face data provides this setting. It allows us to \emph{meta-evaluate} metrics by measuring whether they preserve the ground-truth ordering of identity differences when mesh representation, sampling, support, or pose change. Our experiments show that this problem is not only theoretical. Using a controlled benchmark of synthetic faces derived from a 3D Morphable Model, called REMESH, we observe that standard geometric distances behave inconsistently across mesh types. Raw Chamfer distance is sensitive to topology mismatch. Rigid and non-rigid registration often worsen the agreement with ground-truth identity rankings. We find that this degradation is partly caused by distance compression: registration reduces the spread of inter-subject distances, making different identities appear more similar than they should.
%This distinction matters because, in general, the geometric error between two arbitrary 3D face meshes is only an estimate~\cite{10448898}. A well-defined correspondence-based error exists only under restrictive conditions: the meshes must share a common topology and discretization, and their canonical alignment must be fixed. Under these conditions, one can compute a reference distance directly from corresponding vertices, for example using point-wise RMSE. Synthetic or topology-consistent 3D face data therefore provides a principled way to \emph{meta-evaluate} evaluation metrics: a candidate metric can be judged by how well it preserves the ground-truth ordering of identity-level geometric dissimilarities under changes in mesh representation, sampling, support, and pose. Our experiments show that this distinction between pairwise fit and global ranking fidelity is not merely theoretical. On a controlled benchmark of synthetic 3D faces derived from a 3D Morphable Model, named REMESH, raw Chamfer distance turns out to be strongly affected by topology mismatch, while rigid and non-rigid registration often further reduce agreement with the ground-truth identity ranking. We will show that this degradation is partly explained by distance compression: registration-based pipelines reduce the spread of inter-subject distances, making distinct identities appear artificially more similar.

To address this issue, we define 3D face evaluation as a ranking-preservation problem. Instead of focusing on minimizing surface error for each mesh pair, we ask whether distances are comparable across many subjects. We propose an encoder-only, mesh-agnostic learned metric that maps heterogeneous 3D face meshes into a shared latent space. Pairwise distance is computed in this space using Euclidean distance. The model is trained on REMESH using correspondence-based distances as supervision, with objectives that preserve metric structure, ranking, and identity consistency. At inference time, each mesh is encoded independently. No pairwise ICP, non-rigid registration, landmarks, or dense correspondence estimation are required. Our results show that the learned latent metric better preserves ground-truth identity rankings across topology changes and perturbations. The main contributions are:
%To address the above issue, we formulate robust 3D face evaluation as a ranking-preservation problem. Rather than asking whether an alignment pipeline can minimize local surface discrepancy for each mesh pair, we ask whether the resulting distances remain globally comparable across subject pairs. To this aim, we propose an encoder-only, mesh-agnostic learned metric that maps heterogeneous 3D face meshes into a shared latent space. Pairwise dissimilarity is then computed as Euclidean distance between embeddings. The model is trained on the REMESH topology-consistent data using correspondence-based distances as ground-truth, with objectives that preserve metric stress, pairwise rankings, and identity consistency across mesh realizations. At inference time, each mesh is encoded independently, without pairwise ICP, non-rigid registration, landmarks, or dense correspondence estimation. We will finally show that such learned latent metric better preserves ground-truth rankings across topology changes and perturbations. Ultimately, our contributions are:
\begin{itemize}
\item We formulate 3D face evaluation as a metric meta-evaluation problem, where distances
are judged by identity-ranking preservation under topology, sampling, support,
and pose changes.
\item We propose an encoder-only, mesh-agnostic learned metric that compares
heterogeneous 3D meshes without pairwise registration or dense correspondence at
inference time.
\item We introduce the REMESH controlled cross-topology evaluation protocol and
show that rigid and non-rigid registration can degrade identity-ranking fidelity
by compressing inter-subject distances.
\item We release code, checkpoints, reproducibility artifacts, subject splits,
ground-truth distance matrices, and benchmark metadata; meshes are not
redistributed due to 3DMM licensing constraints.
\end{itemize}
%\begin{itemize}
%    \item We formulate 3D face evaluation as a metric meta-evaluation problem, where candidate distances are assessed by their ability to preserve correspondence-based ground-truth identity rankings under topology, sampling, support, and pose changes.

%    \item We introduce an encoder-only, mesh-agnostic learned metric that embeds heterogeneous 3D face meshes into a shared latent space and compares them without pairwise registration or dense correspondence at inference time.

%    \item We construct and evaluate a controlled REMESH cross-topology protocol with multiple mesh realizations per identity, enabling systematic analysis of topology-induced failures in geometric metrics.

%    \item We show that rigid and non-rigid registration can degrade identity-ranking fidelity by compressing inter-subject distance distributions, despite improving pairwise surface fit.

%    \item We publicly release the trained model checkpoint, training/evaluation code, eproducibility artifacts, subject splits, ground-truth distance matrices, and benchmark metadata. Due to licensing constraints of the underlying 3D morphable model, we do not redistribute the generated meshes directly; instead, we will provide the parameters and metadata required to reproduce the benchmark under the appropriate license.
%\end{itemize}


\section{Related Work}
\label{sec:related_work}
\paragraph{3D face reconstruction evaluation.}
Several solutions have been proposed to evaluate 3D face reconstruction accuracy, most of which work by computing the average point-to-point nearest-neighbor distance after either rigid registration via ICP~\cite{jackson17,deng2019accurate,gecer2019ganfit,gecer2021fast,genova18,Bai_2021_CVPR, prnet, feng2018evaluation, sariyanidi20, shang20} or landmarks~\cite{liu2018disentangling,zhu2020beyond}, or non-rigid alternatives~\cite{ferrari2021sparse,booth20163d,gilani2017dense}.
NoW~\cite{now} evaluates reconstructed meshes by applying landmark-initialized rigid alignment to ground-truth scans and then computing scan-to-mesh distances. This makes the protocol applicable to methods with different output templates, but the resulting error depends on rigid registration quality and closest-point surface matching. 
REALY~\cite{REALY}, instead, uses a structured, region-aware non-rigid registration protocol to obtain localized reconstruction errors. While this provides a more detailed geometric assessment, it requires stronger assumptions on topology, region definitions, and preprocessing. These benchmarks standardize evaluation, but they treat registration-based geometric error as the main proxy for reconstruction quality.

\paragraph{Metric meta-evaluation.}
A complementary line of work asks whether the metrics used by benchmarks are themselves reliable. Synthetic or topology-consistent 3D face data makes this possible because dense correspondences are known, allowing estimated errors to be compared against correspondence-based reference distances. Sariyanidi et al.~\cite{10448898} and M3DFB~\cite{fg2025} followed this meta-evaluation perspective, showing that different alignment and registration pipelines can produce substantially different estimates of reconstruction error, making such assessment unreliable in practice. Our work builds on this idea, but shifts the focus from estimating local geometric error to preserving global identity-level rankings across heterogeneous mesh representations.

\paragraph{Geometric and learned distances.}
Geometric distances such as point-to-point, point-to-surface, and Chamfer distance are simple and widely used, but they are sensitive to sampling density, partial support, topology changes, and residual misalignment~\cite{achlioptas2018learning,nguyen2021point}. Other distances were investigated that are irrespective of mesh parameterization or rotations to eliminate the need for alignment, such as the varifold metric~\cite{kaltenmark2017general, pierson20223d}. However, a kernel needs to be defined to induce the metric on the shapes, which might have varying properties. Learned distances for point clouds, such as DPDist~\cite{urbach2020dpdist}, offer a data-driven alternative to hand-designed geometric metrics, while feature-space metrics in 2D vision, such as LPIPS~\cite{lpips}, show that learned representations can better capture task-relevant similarity than raw signal-space errors. However, learned metrics also depend on their training distribution and are rarely evaluated under controlled topology shifts with known 3D correspondences. In contrast, we train and evaluate a mesh-agnostic learned metric specifically for ranking-preserving 3D face evaluation, and explicitly compare it against registration-based geometric pipelines.

\section{Problem Formulation}
\label{sec:problem_formulation}
We consider the problem of evaluating a distance function between 3D faces.
Given two observed meshes $X_i$ and $X_j$, a candidate metric produces an estimated dissimilarity
%
%\begin{equation}
$D_{\mathrm{ES}}(X_i,X_j) \in \mathbb{R}_{\ge 0}$.
%\label{eq:estimated_dissimilarity}
%\end{equation}

For standard geometric pipelines, $D_{\mathrm{ES}}$ may correspond to Chamfer distance, rigidly aligned point-to-point distance, or non-rigidly registered point-to-surface distance. For the proposed method, $D_{\mathrm{ES}}$ is the Euclidean distance between learned latent embeddings.

In general, the true geometric discrepancy between two arbitrary 3D face meshes is not directly observable. As the meshes may differ in topology, resolution, sampling density, surface support, or pose, the correct point correspondence is unknown. We therefore follow a metric meta-evaluation setting. Starting from topology-consistent synthetic identities with known dense correspondence, we define a reference correspondence-based identity dissimilarity
%
\begin{equation}
D_{\mathrm{GT}}(i,j)
=
\sqrt{
\frac{1}{N}
\sum_{k=1}^{N}
\left\|
v_{i,k} - v_{j,k}
\right\|_2^2
}, 
\label{eq:ground_truth_distance}
\end{equation}

where $v_{i,k}$ and $v_{j,k}$ are corresponding vertices of index $k$ of subjects $i$ and $j$. This distance is used only as a reference for supervision and meta-evaluation; it is not assumed to be available in practice.

Given a set of held-out subjects $\mathcal S$, we compare a candidate distance matrix $D_{\mathrm{ES}}$ against the reference matrix $D_{\mathrm{GT}}$ over all cross-subject pairs. Our primary criterion is Spearman rank correlation,
\begin{equation}
\rho_{\mathrm{S}}
=
\operatorname{corr}_{\mathrm{Spearman}}
\left(
\{D_{\mathrm{ES}}(i,j)\}_{i<j},
\{D_{\mathrm{GT}}(i,j)\}_{i<j}
\right), 
\label{eq:spearman_correlation}
\end{equation}
which measures whether the estimated metric preserves the ordering induced by the ground-truth identity dissimilarities. This is the central requirement in our setting: if subject $i$ is closer to subject $j$ than to subject $k$ under the canonical correspondence-based geometry, a reliable metric should preserve this ordering even when the observed meshes differ in topology, sampling, or pose.

As a secondary criterion, we report Pearson correlation,
\begin{equation}
\rho_{\mathrm{P}}
=
\operatorname{corr}_{\mathrm{Pearson}}
\left(
\{D_{\mathrm{ES}}(i,j)\}_{i<j},
\{D_{\mathrm{GT}}(i,j)\}_{i<j}
\right), 
\label{eq:pearson_correlation}
\end{equation}
which measures linear agreement and therefore captures whether relative distance proportions are preserved beyond ranking. Throughout the paper, higher correlation indicates better agreement with the correspondence-based reference.

\section{Method}
\label{sec:method}

\subsection{Mesh-Agnostic Encoder}
\label{sec:encoder}
We learn a shared embedding space for heterogeneous 3D face meshes. By mesh-agnostic, we mean that the encoder does not require a fixed vertex count, mesh connectivity, or template correspondence across inputs. Given an input mesh $M$, an encoder $f_\theta$ maps it to a latent vector $z=f_\theta(M)\in\mathbb{R}^d$. The learned dissimilarity between two meshes is defined as Euclidean distance in latent space:
\begin{equation}
D_{\theta}(M_i,M_j)
=
\left\|
f_{\theta}(M_i)-f_{\theta}(M_j)
\right\|_2 .
\end{equation}
Unlike registration-based geometric pipelines, this distance does not require pairwise alignment, ICP, non-rigid registration, landmarks, or dense correspondence at inference time.

We use an encoder-only surface network based on DiffusionNet~\cite{diffusionet}. 
For each triangular mesh $M=(V,F)$, vertices are centered and normalized, and the network receives normalized XYZ coordinates together with precomputed geometric operators, including mass, Laplacian, eigensystem, and gradient operators. These operators provide a geometry-aware bias that helps generalization across mesh connectivity and sampling changes. Since the encoder also receives XYZ coordinates, it is not guaranteed to be exactly invariant to rigid transformations; robustness to rotation and translation is instead encouraged by cross-topology training and online perturbation augmentation.

The reported model uses four DiffusionNet blocks with hidden width 128 and outputs a 256-dimensional latent embedding. Vertex features are aggregated by concatenated mean and max pooling followed by a linear projection. The same encoder is shared across all topology variants.
\subsection{Training Objective}
\label{sec:training_objective}

Training uses topology-consistent synthetic identities with an exact subject-level
ground-truth distance matrix $D_{\mathrm{GT}}$. Each mini-batch contains multiple
subjects and multiple mesh realizations per subject. Let
$z_{s,m}=f_\theta(x_{s,m})$ be the embedding of realization $m$ of subject $s$,
and define the batch subject embedding as
$\bar z_s = M_s^{-1}\sum_{m=1}^{M_s} z_{s,m}$.

The model is trained with a mixed subject- and mesh-level objective:
\begin{equation}
\mathcal L =
\lambda_{\mathrm{subj}}
\left(
\mathcal L^{\mathrm{subj}}_{\mathrm{stress}}
+
\lambda_r \mathcal L^{\mathrm{subj}}_{\mathrm{rank}}
\right)
+
\lambda_{\mathrm{mesh}}
\left(
\mathcal L^{\mathrm{mesh}}_{\mathrm{stress}}
+
\lambda_r \mathcal L^{\mathrm{mesh}}_{\mathrm{rank}}
\right)
+
\lambda_i \mathcal L_{\mathrm{id}} .
\end{equation}
Subject-level terms preserve the identity manifold defined by $D_{\mathrm{GT}}$,
while mesh-level terms expose the model to topology-induced variation. We use
$\lambda_{\mathrm{subj}}=\lambda_{\mathrm{mesh}}=1$, $\lambda_r=0.5$, and
$\lambda_i=0.25$.

\paragraph{Stress and ranking preservation.}
For embeddings $Z$ and valid pairs $\Omega$, latent and ground-truth distances
are matched with a scale-normalized Smooth L1 stress:
\begin{equation}
\mathcal L_{\mathrm{stress}}
=
\frac{1}{|\Omega|}
\sum_{(i,j)\in\Omega}
\operatorname{SmoothL1}
\left(
\frac{\lVert z_i-z_j\rVert_2}{\mu_z+\epsilon},
\frac{D_{\mathrm{GT},ij}}{\mu_{\mathrm{GT}}+\epsilon}
\right),
\end{equation}
where $\mu_z$ and $\mu_{\mathrm{GT}}$ are mean distances over $\Omega$. We apply
this loss at both subject and mesh level. We additionally enforce the ordering
induced by $D_{\mathrm{GT}}$ using a margin-based ranking loss. For subject
triplets $(i,j,k)$ satisfying
$D_{\mathrm{GT}}(i,j)>D_{\mathrm{GT}}(i,k)+\tau$,
\begin{equation}
\mathcal L^{\mathrm{subj}}_{\mathrm{rank}}
=
\mathbb E
\left[
\max
\left(
0,\;
\gamma -
\left(
\lVert \bar z_i-\bar z_j\rVert_2
-
\lVert \bar z_i-\bar z_k\rVert_2
\right)
\right)
\right].
\end{equation}
The mesh-level ranking term applies the same principle to pairs of mesh
distances, with ground-truth ordering inherited from subject identities. We use
$\gamma=0.05$ and ignore comparisons whose ground-truth distance difference is
below $\tau=0.02$.

\paragraph{Identity consistency.}
Finally, $\mathcal L_{\mathrm{id}}$ penalizes the variance of different
mesh-realization embeddings around their subject mean $\bar z_s$, encouraging
topology robustness while the stress and ranking terms preserve the global
identity-space geometry. The full mesh-level ranking and identity-consistency
definitions are given in Appendix~\ref{app:mesh_level_ranking_loss}.

\section{Experimental Setup}
\label{sec:experimental_setup}

\subsection{REMESH Cross-Topology Benchmark}
\label{sec:remesh_benchmark}
We build REMESH as a controlled metric benchmark from canonical 3D face shapes generated by the Basel Face Model (BFM)~\cite{bfm}, a 3D morphable model. These meshes are not outputs of image-based reconstruction methods. Instead, they provide reference identity geometry with known dense correspondence, allowing evaluation metrics and registration pipelines to be meta-evaluated under topology, support and sampling changes.

Following the metric meta-evaluation perspective of
Sariyanidi et al.~\cite{10448898}, synthetic data is used here as a controlled tool rather than as an end task. Because dense correspondence is known in the canonical representation, we can define reference identity distances and measure whether candidate metrics preserve them when the observed mesh representation is changed. The source meshes are first represented in a common canonical coordinate system. All meshes are centered and rescaled to an approximately $[-1,1]$ coordinate range. The GT-ready dataset contains 4,999 full-resolution meshes, each with approximately 53k vertices and 106k faces. We extract a canonical facial region using a fixed vertex-index mask, yielding meshes with 23,470 vertices and 46,440 faces.

From these canonical meshes, we construct REMESH, a controlled cross-topology benchmark with 500 identities and six mesh realizations per identity, for a total of 3,000 meshes. The six realizations are: \textbf{Original}, the canonical cropped mesh;
\textbf{Noisy}, the original topology with Gaussian vertex noise; 
\textbf{Remesh}, a smoothed and quadric-decimated mesh with approximately 16.5k vertices;
\textbf{Crop}, a conservative boundary crop with approximately 21k vertices;
\textbf{Down8k}, a strongly decimated mesh with approximately 8.1k vertices; and
\textbf{Up60k}, an upsampled mesh with approximately 60.4k vertices.
Each realization preserves the same underlying identity while changing topology, resolution, sampling density, support, or local geometry.

All meshes are additionally normalized per mesh before any distance is computed, by
centering on the vertex mean and dividing by the largest absolute coordinate. This step is
what makes meshes of different scales comparable, but it is not a neutral choice: the
divisor is a single extreme coordinate and is therefore itself topology-dependent. We
quantify its effect, and report all geometric baselines under a second, area-based
normalization, in Supplementary Sec.~1; the ranking conclusions of this paper hold under
both, and the normalization used here is the stronger of the two for the geometric
baselines in aggregate.

Because all identities of REMESH come from a single morphable model, ``cross-topology''
throughout this paper means \emph{across mesh realizations of one model family}. Transfer
to a second model family is a separate question, addressed in
Sect.~\ref{sec:ood_transfer}.

\subsection{Evaluation Protocol and Implementation Details}
\label{sec:evaluation_protocol}
We use a fixed subject-disjoint split: 400 identities are used for training and 100 identities are held out for evaluation. No held-out identity is used during training. Model selection is performed on the fixed held-out REMESH benchmark, and all reported REMESH numbers use this same subject-disjoint benchmark. We therefore refer to these results as held-out benchmark performance rather than as an independent test-set estimate.

For each held-out identity, all six topology variants are evaluated, yielding 600 evaluation meshes. We consider cross-subject comparisons only. For topology-pair analysis, each ordered pair $(\tau_a,\tau_b)$ is evaluated separately and same-topology pairs are excluded. This gives $6\times5=30$ ordered topology-pair settings. Each setting contains $\binom{100}{2}=4{,}950$ subject pairs, for a total of $148{,}500$ cross-topology mesh-pair observations.

Candidate distances are compared against the correspondence-based $D_{\mathrm{GT}}$ using Spearman correlation as the primary metric and Pearson correlation as a secondary metric, as defined in Sect.~\ref{sec:problem_formulation}. We evaluate clean meshes and four perturbation regimes: jitter, translation, rotation, and mixed. Unless otherwise specified, the REMESH perturbation scale is $\sigma=0.02$.

The reported model uses the \texttt{xyz\_dn} DiffusionNet encoder with latent dimension 256, width 128, four residual blocks, dropout 0.1, and mean--max pooling. Training uses mini-batches of 5 subjects with up to 6 realizations per subject. We train for 120 epochs with learning rate $10^{-4}$, weight decay $10^{-6}$, and gradient clipping at 1.0. Online perturbations are applied during training with probability 0.6, using translation, rotation, and jitter sampled with relative weights $4{:}2{:}1$.

\paragraph{Hardware and software environment.}
The main training and evaluation runs were performed on a workstation with an AMD Ryzen 7 9800X3D CPU
(8 cores, 16 threads), 123 GiB of RAM, and one NVIDIA GeForce RTX 5090 GPU
with 32 GiB of memory, using NVIDIA driver 570.211.01. The software environment
used Python 3.12.3, PyTorch 2.11.0, and CUDA 12.9. On this setup, one 120-epoch
training run takes approximately 24 hours, corresponding to about 12 minutes per
epoch. Bootstrap, correlation, and table-generation analyses operate on
precomputed pair-level distances and run on CPU. Exact reproduction commands are
included in the released artifact bundle.

\section{Results}
\label{sec:results}

\subsection{Cross-Topology Ranking}
\label{sec:xtopo_results}
We first evaluate whether candidate distances preserve the ground-truth identity ranking under topology changes. For each ordered topology pair $\tau_a \rightarrow \tau_b$, we compute pairwise distances over all $\binom{100}{2}=4950$ held-out subject pairs and correlate them with $D_{\mathrm{GT}}$ as reported in Table~\ref{tab:clean_xtopo_chamfer_latent_matrices}. Subject-level bootstrap confidence intervals for all topology-pair entries are reported in Appendix~\ref{app:xtopo_bootstrap_ci}.

\begin{table*}[!ht]
\centering
\scriptsize
\caption{Clean same-topology and cross-topology \textbf{Spearman correlation} ($\uparrow$) by ordered topology pair. Rows indicate source topology and columns indicate target topology. Diagonal entries are same-topology evaluations; off-diagonal entries are cross-topology evaluations. Bold indicates the better score between Raw Chamfer and latent distance for the same ordered topology pair.}
\label{tab:clean_xtopo_chamfer_latent_matrices}
\setlength{\tabcolsep}{3.2pt}
\renewcommand{\arraystretch}{1.08}
\resizebox{\textwidth}{!}{%
\begin{tabular}{lcccccc|cccccc}
\toprule
& \multicolumn{6}{c}{\textbf{Raw Chamfer}}
& \multicolumn{6}{c}{\textbf{Latent distance}} \\
\cmidrule(lr){2-7}\cmidrule(lr){8-13}
\textbf{Source}
& crop & down8k & noisy & original & remesh & up60k
& crop & down8k & noisy & original & remesh & up60k \\
\midrule
crop
& 0.558 & -0.093 & -0.048 & -0.035 & -0.008 & -0.074
& \textbf{0.831} & \textbf{0.668} & \textbf{0.737} & \textbf{0.760} & \textbf{0.658} & \textbf{0.687} \\
down8k
& 0.037 & 0.743 & 0.472 & 0.407 & 0.436 & 0.611
& \textbf{0.692} & \textbf{0.883} & \textbf{0.793} & \textbf{0.781} & \textbf{0.795} & \textbf{0.804} \\
noisy
& 0.083 & 0.313 & 0.739 & 0.659 & 0.579 & 0.489
& \textbf{0.747} & \textbf{0.779} & \textbf{0.900} & \textbf{0.821} & \textbf{0.771} & \textbf{0.791} \\
original
& 0.093 & 0.240 & 0.627 & 0.729 & 0.553 & 0.414
& \textbf{0.771} & \textbf{0.772} & \textbf{0.822} & \textbf{0.902} & \textbf{0.762} & \textbf{0.787} \\
remesh
& 0.102 & 0.338 & 0.621 & 0.622 & 0.709 & 0.525
& \textbf{0.685} & \textbf{0.796} & \textbf{0.784} & \textbf{0.771} & \textbf{0.873} & \textbf{0.806} \\
up60k
& 0.061 & 0.564 & 0.597 & 0.539 & 0.573 & 0.735
& \textbf{0.705} & \textbf{0.798} & \textbf{0.798} & \textbf{0.788} & \textbf{0.801} & \textbf{0.890} \\
\bottomrule
\end{tabular}
}
\end{table*}

Raw Chamfer is competitive when source and target meshes share the same or a very similar representation, but it deteriorates sharply under support and topology mismatch. The failure is most visible for crop-involving pairs, where Chamfer correlations often become close to zero or negative. In contrast, the learned metric remains strongly correlated with the ground-truth identity ranking across all topology pairs. The appendix confidence intervals show that this pattern is stable under subject-level bootstrap resampling. This indicates that the gain is not only an aggregate effect, but a systematic reduction of topology-induced evaluation bias.

\subsection{Does Alignment Help?}
\label{sec:alignment_results}

We next compare in Table~\ref{tab:alignment_effect} the learned metric with increasingly stronger geometric alignment pipelines. The goal is to test whether reducing pairwise geometric discrepancy also improves identity-ranking fidelity. We compare raw Chamfer, rigid ICP followed by Chamfer distance, non-rigid ICP point-to-point distance (NICP P2P), and non-rigid ICP point-to-triangle distance (NICP P2Tri). To avoid making the comparison depend only on the most severe crop failures, we also report a no-crop cross-topology setting in which all crop-involving pairs are removed for every method.

We report 95\% subject-level bootstrap confidence intervals using 1,000
replicates. Each replicate samples held-out subjects with replacement, rebuilds
the induced cross-subject pairs, and recomputes Spearman correlation. This avoids
treating mesh pairs sharing the same subject as independent. The intervals
capture uncertainty over the held-out subject population under the fixed
checkpoint and evaluation protocol, but not training or split-selection
variability.

\begin{table}[!ht]
\centering
\small
\caption{Effect of rigid and non-rigid alignment on identity-ranking fidelity.
Each cell reports Spearman correlation ($\uparrow$) with $D_{\mathrm{GT}}$ and
95\% subject-level bootstrap confidence intervals in brackets. The no-crop
setting removes all topology pairs involving crop.}
\label{tab:alignment_effect}
\begin{tabular}{lcc}
\toprule
Method & Original-to-original & No-crop cross-topology \\
\midrule
Latent distance & \textbf{0.902 [0.875, 0.924]} & \textbf{0.868 [0.833, 0.894]} \\
Chamfer & 0.729 [0.667, 0.788] & 0.552 [0.488, 0.606] \\
Rigid ICP + Chamfer & 0.600 [0.528, 0.663] & 0.514 [0.449, 0.573] \\
Rigid ICP + NICP + P2P & 0.456 [0.357, 0.544] & 0.414 [0.326, 0.499] \\
Rigid ICP + NICP + P2Tri & 0.468 [0.373, 0.552] & 0.433 [0.350, 0.514] \\
\bottomrule
\end{tabular}

\vspace{0.3em}
\footnotesize{ICP here is the scaled rigid ICP of the geometry pipeline. A landmark-based
rigid aligner is considerably stronger on this task and is reported, together with the
remaining M3DFB~\cite{fg2025} estimators, in Supplementary Sec.~2.2; the ordering with
respect to the learned metric is unchanged.}
\end{table}

Stronger alignment does not improve ranking fidelity. Even in the original-to-original setting, where topology is fixed, rigid ICP followed by Chamfer distance reduces Spearman correlation from $0.729$ to $0.600$, and non-rigid ICP further reduces it to approximately $0.46$--$0.47$. The same trend holds in the no-crop cross-topology setting: removing crop pairs improves Chamfer, but registration-based distances still remain below both Chamfer and the learned metric. These results support the claim that pairwise alignment is not a neutral preprocessing step for identity-aware evaluation. Alignment can improve local surface fit while reducing the global comparability of distances across subject pairs.

\subsection{Computational Cost}
\label{sec:computational_cost}

The proposed metric changes the scaling of evaluation. Chamfer, ICP, and NICP
pay geometry-processing or registration costs for every mesh pair,
$O(PC_{\mathrm{geom/reg}})$, whereas our metric pays a reusable per-mesh
operator/embedding cost plus cheap latent comparisons,
$O(NC_{\mathrm{op}}+NC_{\mathrm{enc}}+Pd)$. After embeddings are available,
each comparison is only a Euclidean distance in latent space.

\begin{table}[!ht]
\centering
\scriptsize
\caption{Runtime scaling on a 742,500-pair evaluation block. Latent distance
excludes one-time preprocessing; full FaceBench ICP/NICP timings are
order-of-magnitude extrapolations from subset runs.}
\label{tab:runtime_scaling}
\begin{tabular}{lc}
\toprule
Method & Estimated time \\
\midrule
Latent distance & $<$1 min \\
Optimized NICP correspondence & $\sim$2.0 h \\
Rigid ICP + Chamfer & $\sim$3.1 h \\
Raw Chamfer & $\sim$6.1 h \\
Full FaceBench rigid ICP & $\sim$5.8 days$^\dagger$ \\
Full FaceBench NICP & $\sim$21 days$^\dagger$ \\
\bottomrule
\end{tabular}

\vspace{0.3em}
\footnotesize{$^\dagger$ Extrapolated from subset runs with the full FaceBench pipeline.}
\end{table}

In our implementation, constructing DiffusionNet operators with
$k_{\mathrm{eig}}=128$ takes 1.65 s at 8.1k vertices, 3.70 s at 16.5k, 5.71 s at
23.5k and 15.65 s at 60.4k, measured on a dedicated node as the median of three
trials per mesh. The cost is essentially serial: repeating the 23.5k-vertex mesh
at 1, 4, 8, 16 and 32 threads gives 5.82, 6.22, 5.44, 5.31 and 5.32 s, since the
sparse Lanczos iteration that dominates offers little to parallelize. Encoding is
then performed once per mesh and costs a further 0.73 s on average, so a complete
pipeline is about 6.8 s per mesh. For the full REMESH benchmark of 3,000 meshes
this is a one-time cost of roughly five core-hours, or tens of minutes with modest
parallelization, together with 20.4 MB of stored operators per mesh. The expensive
part therefore scales linearly in the number of meshes and is amortized across all
$O(N^2)$ pairwise comparisons; on the 742,500-pair block of
Table~\ref{tab:runtime_scaling} it is a few percent of the alternatives' cost.

These figures replace lower ones quoted in an earlier version of this work, which
understated the operator cost by roughly a factor of three. The amortization
argument is unaffected, since the cost is incurred per mesh rather than per pair,
but the numbers themselves were wrong and are corrected here.

\subsection{Why Alignment Hurts: Distance Compression}
\label{sec:distance_compression}
The degradation above is not a failure of ICP as a pairwise fitting algorithm,
but a mismatch between pairwise registration and identity-aware evaluation. For a
mesh pair $(M_i,M_j)$, registration estimates
\begin{equation}
T_{ij}^{*}=\arg\min_T d_{\mathrm{geom}}(T(M_i),M_j),
\end{equation}
so each pair is evaluated after a different transformation chosen to minimize
local discrepancy. For 3D faces, this can absorb identity-related morphology
such as facial depth, jaw profile, forehead shape, and local protrusions, making
different identities appear artificially more similar
(Fig.~\ref{fig:icp_identity_suppression}).

\begin{figure}[!ht]
    \centering
    \includegraphics[width=0.95\linewidth]{imgs/rigid_mixed_sigma0.20_id0084_GTready_down8k__to__id0478_GTready_up60k_paper_identity.png}
    \caption{
    Qualitative illustration of identity-signal suppression by pairwise ICP.
    The figure compares two distinct identities before and after pairwise
    alignment. In the common reference space, morphology differences such as
    profile depth, forehead shape, and lip protrusion remain visible. After
    pairwise ICP, the source is transformed to improve local surface overlap with
    the target, partially absorbing identity variations.
    Pair-specific transformations reduce comparability across subject
    pairs.
    }
    \label{fig:icp_identity_suppression}
\end{figure}

We quantify this effect over 148{,}500 clean cross-topology mesh pairs by
measuring the spread of inter-subject distance distributions. If registration
preserves identity-relevant variation, the distribution should retain dynamic
range; if it suppresses such variation, the distribution should compress. We
report the interquartile range (IQR), standard deviation, range, and normalized
spread ratios relative to raw Chamfer in Table~\ref{tab:distance_compression}.

\begin{table}[!ht]
\centering
\small
\caption{
Distance distribution compression on clean cross-topology mesh pairs.
IQR denotes the interquartile range, $Q_{75}-Q_{25}$, and Std. denotes the
standard deviation of the inter-subject distance distribution. IQR/raw and
Std./raw are normalized spread ratios computed with respect to raw Chamfer.
Lower normalized ratios indicate stronger compression of inter-subject distances.
}

\label{tab:distance_compression}
\begin{tabular}{lcccccc}
\toprule
Metric & Median & IQR & Std. & Range & IQR/raw & Std./raw \\
\midrule
Raw Chamfer & 0.0445 & 0.0202 & 0.0166 & 0.1133 & 1.000 & 1.000 \\
Rigid ICP P2P & 0.0351 & 0.0108 & 0.0092 & 0.0790 & 0.537 & 0.557 \\
NICP P2P & 0.0288 & 0.0044 & 0.0036 & 0.0344 & 0.219 & 0.216 \\
NICP P2Tri & 0.0258 & 0.0048 & 0.0038 & 0.0360 & 0.236 & 0.232 \\
\bottomrule
\end{tabular}
\end{table}

Registration substantially compresses inter-subject distances. Rigid ICP retains
only 53.7\% of the raw Chamfer IQR, while NICP P2P retains only 21.9\%. Thus,
subject pairs that should remain separated under the ground-truth ranking are
mapped to a much narrower registered-distance range
(Fig.~\ref{fig:registration_distance_compression}).

\begin{figure}[!ht]
    \centering
    \includegraphics[width=0.9\linewidth]{imgs/raw_vs_registered_density_panels_shared_axes.png}
    \caption{Registration compresses the distance landscape. Each panel compares raw Chamfer distances with distances obtained after a registration-based pipeline on clean cross-topology mesh pairs. The diagonal indicates identity mapping, while the binned median shows that non-rigid ICP maps pairs with substantially different raw distances into a narrow registered-distance range.}
\label{fig:registration_distance_compression}
\end{figure}

This compression explains why stronger geometric fitting can reduce agreement with ground-truth identity rankings. Registration does not merely rescale distances; it reduces the dynamic range needed to distinguish subject pairs. The learned metric avoids this pair-specific fitting step by embedding all meshes into a shared space and comparing them with a single global distance function.

\subsection{Out-of-Distribution Transfer}
\label{sec:ood_transfer}
We finally evaluate transfer to FaceScape~\cite{facescape}, a separate mesh family unseen during REMESH training. The goal is not to claim universal zero-shot generalization, but to assess how much of the learned identity-ranking structure transfers beyond the source benchmark.

As shown in Table~\ref{tab:faceverse_summary}, in the full-neutral FaceScape setting, the REMESH-trained model shows nontrivial zero-shot transfer: on 110 neutral-pose subjects and 5,995 subject pairs, the latent distance reaches Spearman correlation $0.532$ in the clean setting. However, Chamfer remains stronger, reaching $0.605$. Under mixed perturbations, the gap increases, with latent Spearman $0.477$ compared with $0.638$ for Chamfer. Thus, the learned metric does not collapse out of domain, but it does not outperform Chamfer in this within-topology FaceScape setting.

The harder FaceScape cross-topology setting is where the first version of this work
reported a clear limitation: representing each subject by a downsampled FaceScape mesh and
an independently remeshed 10k-vertex variant, the latent metric reached only $0.115$
Spearman against $0.598$ for Chamfer, which we read as a substantial cross-model domain
shift. A controlled decomposition shows that this reading was wrong, and that the
experiment confounded two changes: the morphable model family, and the rule used to build
the second topology of each pair.

The 10k-vertex mate in that experiment is a Poisson surface reconstruction, which closes
the boundary loops of its source, adds $\approx21\%$ surface area, leaves the original
support, and departs from its own source surface by a normalized Hausdorff distance of
$0.38$ --- against $0.088$ for the \texttt{original}/\texttt{remesh} pair the metric was
trained on. Rebuilding the mate with the REMESH \texttt{remesh} rule, and changing nothing
else, raises FaceScape cross-topology transfer from $0.057$ to $\mathbf{0.404}$; our
synthetic FLAME half of the benchmark reaches $0.406$ under the same rule, i.e. a second
model family costs almost nothing. Conversely, applying the Poisson construction to FLAME
drops it from $0.406$ to $0.109$, reproducing the apparent cross-model failure
\emph{without changing the morphable model at all}. The weakness is therefore robustness to
surface re-reconstruction; it is present in every domain including the training one, and it
is a different axis from cross-model transfer. Both axes are reported in Supplementary
Sec.~4, together with the observation that matching the mesh \emph{support} to the training
domain lowers rather than raises the correlation, because $D_{\mathrm{GT}}$ integrates over
facial regions a crop discards.

Few-shot FaceScape adaptation substantially improves this behavior. With 75 FaceScape fine-tuning identities and evaluation on 35 held-out subjects, the adapted model reaches $0.537$ Spearman in the clean cross-topology setting, outperforming Chamfer at $0.403$. It also remains stronger under translation with $\sigma=0.10$ ($0.486$ vs. $0.417$). These results suggest that the FaceScape gap is primarily a domain-shift issue rather than a fundamental failure of the learned metric, although robustness under stronger jitter remains limited. Full FaceScape adaptation and forgetting results are reported in
Appendix~\ref{app:faceverse_adaptation}.

\begin{table}[!ht]
\centering
\small
\caption{Zero-shot transfer on FaceScape is nontrivial in the within-topology full-neutral setting, insufficient in the harder cross-topology setting, and improves after few-shot target-domain adaptation.}
\label{tab:faceverse_summary}
\begin{tabular}{lccc}
\toprule
Setting & Latent & Chamfer & $\Delta$ \\
\midrule
Full-neutral zero-shot, clean & 0.532 & 0.605 & -0.073 \\
Full-neutral zero-shot, mixed & 0.477 & 0.638 & -0.161 \\
Cross-topology zero-shot, Poisson mate & 0.057 & 0.598 & -0.541 \\
Cross-topology zero-shot, \textsc{remesh} mate & 0.404 & --- & --- \\
75-shot cross-topology, clean & \textbf{0.537} & 0.403 & +0.134 \\
75-shot cross-topology, trans. 0.10 & \textbf{0.486} & 0.417 & +0.069 \\
\bottomrule
\end{tabular}

\vspace{0.3em}
\footnotesize{The two cross-topology rows differ only in how the second topology of each
pair is constructed; see Sect.~\ref{sec:ood_transfer}. The $0.057$ figure is our
re-measurement of the first version's $0.115$ under the evaluator used for all rows here.}
\end{table}


\section{Dataset and Code Availability}
\label{sec:availability}

The anonymized code repository is available at:
\url{https://anonymous.4open.science/r/WBES-FaceEmbedding-120B/README.md}.
The repository includes training and inference scripts, evaluation code,
configuration files, and a reproducibility README describing how to reproduce the
main tables and figures. The full non-anonymized code repository, trained model checkpoint, model card,
and reproducibility artifact bundle will be released publicly upon acceptance.
%The artifact bundle will include trained checkpoints, evaluation outputs,
%generated tables, plots, configuration files and auxiliary metadata. 
%used to produce the results reported in the paper. This is intended to make the numerical results auditable and to reduce the computation required to inspect or reproduce the reported comparisons.

%The REMESH benchmark is intended to support metric meta-evaluation rather than
%to provide a new real-world 3D face reconstruction dataset. The synthetic face
%geometry is generated from a 3D morphable model. 
Because the underlying 3DMM used to generate the synthetic REMESH data
requires a separate license agreement, we do not redistribute the original mesh
assets directly. Instead, we will release the generation coefficients and the
benchmark metadata needed to reconstruct the evaluation protocol. 
%including subject split files, ground-truth identity-distance matrices, topology realization metadata, and evaluation resources. %This design preserves the
%reproducibility of the metric meta-evaluation while respecting the licensing
% constraints of the underlying morphable model.

\section{Conclusions and Limitations}\label{sec:limitations}
We investigated a key limitation of standard error-computation pipelines for
3D face reconstruction. Our experiments show that pre-processing and alignment
steps commonly used in registration-based pipelines can distort the global
structure of shape similarity, making distances less consistent across mesh
pairs and degrading identity-ranking fidelity. This occurs because pair-specific
surface-fit objectives can suppress identity-related morphology and non-uniformly
compress inter-subject distances. Therefore, reconstruction evaluation should
distinguish global canonical normalization from pair-specific alignment, and
should not assume that lower post-registration surface error implies better
identity-aware comparability.

We proposed an encoder-only, mesh-agnostic learned metric trained to make
latent-space distances follow correspondence-based shape dissimilarities. The
metric is designed for ranking-preserving evaluation, not calibrated physical
measurement: latent distances should be interpreted as relative dissimilarities,
not as millimeter- or centimeter-level errors. As with all learned metrics,
performance depends on the training distribution --- but our controlled decomposition of
the FaceScape results (Sect.~\ref{sec:ood_transfer}) locates that dependence more
precisely than a generic domain-shift statement would. Under a matched pair construction
the metric transfers to a second model family, including real scans, at essentially the
same level as to a second synthetic one; what it is not robust to is a mate produced by
surface re-reconstruction, and that weakness is present in the training domain as well.
Multi-domain training and augmentation with re-reconstructed realizations therefore target
the measured failure directly, and are the natural next step, together with calibration
and interpretability of the facial regions driving latent distances.

Three further scoping limits should be stated plainly. The benchmark contains only
\emph{neutral} expressions, so nothing here establishes behaviour under expression
variation. The identities are sampled from a morphable model and therefore inherit the
demographic composition of the data that model was fitted to; a metric trained on them
inherits it in turn, and we make no claim of demographic uniformity. Finally, the
DiffusionNet backbone is discretization-agnostic but \emph{not} invariant to rigid motion;
robustness to pose in our results comes from the training perturbations, not from an
architectural guarantee.

Although the method is intended for evaluation rather than identification, it
operates on facial geometry and could indirectly support identity-related
applications if misused. We therefore frame the release as a benchmark and
metric-evaluation tool, report its domain limitations, and avoid redistributing
licensed mesh assets directly.
{\small
\bibliographystyle{IEEEtran}
\bibliography{references}
}

\appendix

\section{Additional Dataset and Implementation Details}
\label{app:implementation_details}

\subsection{REMESH Construction Details}
\label{app:remesh_details}

The main REMESH construction protocol is described in
Sec.~\ref{sec:remesh_benchmark}. Here we provide additional implementation
details and a qualitative visualization of the resulting topology variants.

The GT-ready source dataset contains 4,999 full-resolution synthetic meshes, each
with 53,215 vertices and 105,954 faces. After canonical facial-region extraction,
each cropped canonical mesh contains 23,470 vertices and 46,440 faces. The
generated REMESH variants have the following approximate resolutions: Remesh,
16.5k vertices and 32.5k faces; Crop, 21k vertices and 41.5k faces; Down8k,
8.1k vertices and 16k faces; and Up60k, 60.4k vertices and 120k faces.

Figure~\ref{fig:remesh_synthetic_examples_app} shows representative REMESH
examples.

\begin{figure}[h]
    \centering
    \includegraphics[width=\linewidth]{imgs/remesh_id0100_six_topologies_pointcloud_colored_rot180_2rows.png}
    \caption{
    REMESH synthetic benchmark examples. Each row shows one synthetic identity
    represented by multiple mesh realizations: original, noisy, remeshed,
    cropped, downsampled, and upsampled. The benchmark evaluates whether a
    distance metric preserves identity-level geometric relationships despite
    changes in mesh representation.
    }
    \label{fig:remesh_synthetic_examples_app}
\end{figure}

\subsection{DiffusionNet Operators}
\label{app:diffusionnet_operators}
For every mesh realization, we precompute the surface operators required by DiffusionNet, including the mass matrix, Laplace--Beltrami operator, spectral eigensystem, and tangent gradient operators. Vertex coordinates encode the extrinsic position of the surface in $\mathbb{R}^3$, whereas the Laplace--Beltrami operators encode intrinsic surface geometry.

Under a rigid motion $T(x)=Rx+t$, with $R\in SO(3)$ and $t\in\mathbb{R}^3$, intrinsic distances, angles, and area elements on the continuous surface are preserved. Consequently, the continuous Laplace--Beltrami operator and its spectrum are invariant to global rotations and translations:
\begin{equation}
\Delta_{T(\mathcal{M})}=\Delta_{\mathcal{M}},
\qquad
\Delta_{\mathcal{M}}\phi_\ell=\lambda_\ell\phi_\ell.
\end{equation}
This provides a useful geometric bias for robustness to global pose changes.

However, topology changes still affect the discrete operators. Remeshing,
decimation, subdivision, cropping, and noise modify the sampled mesh and
therefore change the discrete mass matrix, Laplacian, eigenbasis, and gradient
operators. Thus, intrinsic operators do not guarantee topology invariance by
themselves. Our model combines extrinsic XYZ coordinates with geometry-aware
diffusion operators and is explicitly trained across topology variants, so that
topology-dependent operator differences are treated as nuisance variation while
identity-related geometric differences are preserved.

\subsection{Training Details}
\label{app:training_details}

The reported model uses the \texttt{xyz\_dn} encoder with latent dimension 256,
hidden width 128, four DiffusionNet blocks, dropout 0.1, and mean--max pooling.
Training uses mini-batches of 5 subjects with up to 6 mesh realizations per
subject. The model is trained for 120 epochs with learning rate $10^{-4}$,
weight decay $10^{-6}$, and gradient clipping at 1.0.

The training objective is defined in Sec.~\ref{sec:training_objective}. In all
reported experiments we use $\lambda_{\mathrm{subj}}=1$,
$\lambda_{\mathrm{mesh}}=1$, $\lambda_r=0.5$, and $\lambda_i=0.25$.

Mesh-level supervision uses cross-topology pairs only. Ranking losses sample
1024 constraints per batch, with margin $\gamma=0.05$, tolerance $\tau=0.02$,
and 70\% hard constraints.

During training, online perturbations are applied with probability 0.6. The
perturbation modes are translation, rotation, and jitter, sampled with relative
weights $4{:}2{:}1$. Jitter magnitude is sampled log-uniformly between
$5\times10^{-4}$ and $2\times10^{-2}$. Rotations use angles between
$0.5^\circ$ and $12^\circ$, and translations use per-axis standard deviation
between $0.001$ and $0.03$.
\subsection{Mesh-Level Ranking Loss}
\label{app:mesh_level_ranking_loss}

At the mesh level, we compare two valid mesh pairs $a$ and $b$ and enforce the
ordering induced by their corresponding subject-level ground-truth distances:
\begin{equation}
\mathcal L^{\mathrm{mesh}}_{\mathrm{rank}}
=
\mathbb E_{(a,b)}
\left[
\max
\left(
0,\;
\gamma -
\operatorname{sign}
\left(
D_{\mathrm{GT}}^a-D_{\mathrm{GT}}^b
\right)
\left(
d_z^a-d_z^b
\right)
\right)
\right],
\end{equation}
where $d_z^a$ and $d_z^b$ are the latent distances of the two mesh pairs. We
ignore comparisons whose ground-truth distance difference is below $\tau=0.02$.

The identity-consistency term used in the main objective is
\begin{equation}
\mathcal L_{\mathrm{id}}
=
\frac{1}{|\mathcal S|}
\sum_{s\in\mathcal S}
\frac{1}{M_s}
\sum_{m=1}^{M_s}
\left\|
z_{s,m}-\bar z_s
\right\|_2^2 .
\end{equation}

\section{Cross-Topology Bootstrap Confidence Intervals}
\label{app:xtopo_bootstrap_ci}

Tables~\ref{tab:xtopo_chamfer_ci_app} and~\ref{tab:xtopo_latent_ci_app}
report 95\% subject-level bootstrap confidence intervals for the clean
same-topology and cross-topology ranking experiment. Each cell reports Spearman
correlation with $D_{\mathrm{GT}}$, followed by the 95\% confidence interval in
brackets. Confidence intervals are computed over 1,000 subject-level bootstrap
replicates, using the same procedure described in Sec.~\ref{sec:alignment_results}.

\begin{table*}[!htbp]
\centering
\scriptsize
\caption{
Raw Chamfer clean same-topology and cross-topology Spearman correlation with
95\% subject-level bootstrap confidence intervals. Rows indicate source topology
and columns indicate target topology.
}
\label{tab:xtopo_chamfer_ci_app}
\setlength{\tabcolsep}{2.0pt}
\renewcommand{\arraystretch}{1.08}
\resizebox{\textwidth}{!}{%
\begin{tabular}{lcccccc}
\toprule
Source & crop & down8k & noisy & original & remesh & up60k \\
\midrule
crop & 0.558 [0.480, 0.632] & -0.093 [-0.212, 0.034] & -0.048 [-0.165, 0.072] & -0.035 [-0.146, 0.091] & -0.008 [-0.132, 0.112] & -0.074 [-0.198, 0.049] \\
down8k & 0.037 [-0.095, 0.174] & 0.743 [0.687, 0.792] & 0.472 [0.393, 0.551] & 0.407 [0.310, 0.491] & 0.436 [0.344, 0.519] & 0.611 [0.528, 0.682] \\
noisy & 0.083 [-0.030, 0.204] & 0.313 [0.216, 0.401] & 0.739 [0.674, 0.794] & 0.659 [0.586, 0.724] & 0.579 [0.498, 0.649] & 0.489 [0.406, 0.567] \\
original & 0.093 [-0.014, 0.211] & 0.240 [0.154, 0.328] & 0.627 [0.551, 0.698] & 0.729 [0.667, 0.790] & 0.553 [0.472, 0.624] & 0.414 [0.327, 0.496] \\
remesh & 0.102 [-0.005, 0.218] & 0.338 [0.255, 0.431] & 0.621 [0.541, 0.691] & 0.622 [0.544, 0.691] & 0.709 [0.647, 0.763] & 0.525 [0.444, 0.605] \\
up60k & 0.061 [-0.067, 0.184] & 0.564 [0.479, 0.644] & 0.597 [0.518, 0.669] & 0.539 [0.463, 0.608] & 0.573 [0.497, 0.643] & 0.735 [0.674, 0.791] \\
\bottomrule
\end{tabular}
}
\end{table*}

\begin{table*}[!htbp]
\centering
\scriptsize
\caption{
Latent distance clean same-topology and cross-topology Spearman correlation with
95\% subject-level bootstrap confidence intervals. Rows indicate source topology
and columns indicate target topology.
}
\label{tab:xtopo_latent_ci_app}
\setlength{\tabcolsep}{2.0pt}
\renewcommand{\arraystretch}{1.08}
\resizebox{\textwidth}{!}{%
\begin{tabular}{lcccccc}
\toprule
Source & crop & down8k & noisy & original & remesh & up60k \\
\midrule
crop & 0.831 [0.789, 0.866] & 0.668 [0.603, 0.730] & 0.737 [0.672, 0.791] & 0.760 [0.696, 0.811] & 0.658 [0.591, 0.720] & 0.687 [0.615, 0.747] \\
down8k & 0.692 [0.619, 0.751] & 0.883 [0.847, 0.909] & 0.793 [0.741, 0.839] & 0.781 [0.724, 0.829] & 0.795 [0.746, 0.836] & 0.804 [0.746, 0.847] \\
noisy & 0.747 [0.684, 0.802] & 0.779 [0.722, 0.825] & 0.900 [0.874, 0.923] & 0.821 [0.771, 0.864] & 0.771 [0.720, 0.816] & 0.791 [0.737, 0.837] \\
original & 0.771 [0.715, 0.821] & 0.772 [0.717, 0.816] & 0.822 [0.774, 0.865] & 0.902 [0.873, 0.925] & 0.762 [0.712, 0.806] & 0.787 [0.741, 0.829] \\
remesh & 0.685 [0.615, 0.740] & 0.796 [0.742, 0.840] & 0.784 [0.729, 0.833] & 0.771 [0.716, 0.818] & 0.873 [0.838, 0.900] & 0.806 [0.751, 0.848] \\
up60k & 0.705 [0.631, 0.766] & 0.798 [0.744, 0.844] & 0.798 [0.744, 0.843] & 0.788 [0.728, 0.834] & 0.801 [0.755, 0.845] & 0.890 [0.860, 0.917] \\
\bottomrule
\end{tabular}
}
\end{table*}

\section{Additional Perturbation Results}

\label{app:perturbations}


\subsection{Selected Cross-Topology Pairs}
\label{app:selected_topology_pairs}

Table~\ref{tab:xtopo_latent_chamfer_app} reports cross-topology Spearman
correlation across perturbation regimes for representative ordered topology
pairs. Each entry reports latent distance / raw Chamfer. The results confirm the
trend observed in the main paper: Chamfer is highly sensitive to topology and
support mismatch, while the learned metric preserves substantially stronger
agreement with the ground-truth ranking.

\begin{table*}[h]
\centering
\small
\caption{
Cross-topology Spearman correlation across perturbation regimes for selected
ordered topology pairs. Results are computed over 4,950 subject pairs per
topology pair. Each cell reports \textbf{Latent / Chamfer}. Higher is better.
}
\label{tab:xtopo_latent_chamfer_app}
\begin{tabular}{lccccc}
\toprule
Topology Pair & Clean & Jitter & Translation & Rotation & Mixed \\
\midrule
crop $\rightarrow$ down8k
& 0.668 / -0.093
& 0.585 / -0.092
& 0.653 / -0.092
& 0.648 / -0.097
& 0.557 / -0.093 \\

crop $\rightarrow$ remesh
& 0.658 / -0.008
& 0.564 / -0.014
& 0.647 / 0.005
& 0.632 / -0.013
& 0.518 / -0.021 \\

crop $\rightarrow$ up60k
& 0.687 / -0.074
& 0.628 / -0.070
& 0.672 / -0.069
& 0.650 / -0.081
& 0.569 / -0.073 \\

down8k $\rightarrow$ noisy
& 0.793 / 0.472
& 0.692 / 0.495
& 0.781 / 0.453
& 0.769 / 0.465
& 0.681 / 0.461 \\

down8k $\rightarrow$ up60k
& 0.804 / 0.611
& 0.732 / 0.576
& 0.789 / 0.598
& 0.758 / 0.579
& 0.695 / 0.544 \\

original $\rightarrow$ remesh
& 0.762 / 0.553
& 0.674 / 0.549
& 0.753 / 0.546
& 0.737 / 0.545
& 0.640 / 0.497 \\

remesh $\rightarrow$ up60k
& 0.806 / 0.525
& 0.755 / 0.532
& 0.788 / 0.531
& 0.758 / 0.492
& 0.688 / 0.506 \\
\bottomrule
\end{tabular}
\end{table*}

\subsection{Full Registration Pipeline Under Perturbations}
\label{app:facebench_pipeline_perturbations}

We further compare the proposed latent distance against the full FaceBench
geometry pipeline under perturbations. This benchmark includes raw Chamfer,
rigid ICP point-to-point distance, and non-rigid ICP distances computed as
point-to-point and point-to-triangle errors. Unlike the cross-topology breakdown
in the main paper, this experiment evaluates complete registration pipelines
rather than raw Chamfer alone.

\begin{table*}[h]
\centering
\small
\caption{
FaceBench pipeline comparison under perturbations. Each cell reports Spearman /
Pearson correlation with the ground-truth identity distance. Bold indicates the
best score. Higher is better.
}
\label{tab:facebench_pipeline_perturbed_app}
\begin{tabular}{lcccc}
\toprule
Method & Clean & Translation 0.10 & Rotation 0.10 & Jitter 0.10 \\
\midrule
Latent distance
& \textbf{0.786} / \textbf{0.816}
& \textbf{0.770} / \textbf{0.795}
& \textbf{0.765} / \textbf{0.802}
& \textbf{0.562} / \textbf{0.597} \\

Raw Chamfer
& 0.345 / 0.388
& 0.343 / 0.387
& 0.337 / 0.377
& 0.258 / 0.312 \\

Rigid ICP P2P
& 0.383 / 0.468
& 0.372 / 0.465
& 0.379 / 0.464
& -0.028 / 0.004 \\

NICP P2P
& 0.276 / 0.364
& 0.246 / 0.340
& 0.254 / 0.348
& -0.005 / 0.026 \\

NICP P2Tri
& 0.276 / 0.361
& 0.245 / 0.337
& 0.255 / 0.346
& 0.067 / 0.100 \\
\bottomrule
\end{tabular}
\end{table*}

Across clean, translation, rotation, and jitter perturbations, the learned
latent distance remains more correlated with the ground-truth identity distance
than all geometry-only baselines. The gap is particularly pronounced under
jitter, where rigid and non-rigid registration become nearly uncorrelated with
identity ranking, while the latent distance still preserves a Spearman
correlation of $0.562$.

\section{Additional Registration Compression Analysis}
\label{app:compression_analysis}

Figure~\ref{fig:distance_compression_boxplot_app} reports the distribution of
clean cross-topology geometric distances before and after registration. The
boxplot complements the distance-compression statistics reported in the main
paper. Compared with raw Chamfer, rigid ICP and especially non-rigid ICP reduce
the spread of inter-subject distances, supporting the hypothesis that
registration can suppress identity-discriminative variation.

\begin{figure}[h]
    \centering
    \includegraphics[width=0.85\linewidth]{imgs/clean_geometry_distance_boxplot.png}
    \caption{
    Distribution of clean cross-topology geometric distances before and after
    registration. Compared with raw Chamfer, rigid ICP and especially non-rigid
    ICP substantially reduce the spread of inter-subject distances.
    }
    \label{fig:distance_compression_boxplot_app}
\end{figure}

We observe the same trend under translation and rotation perturbations. Rigid ICP
consistently reduces the interquartile range to approximately 59--63\% of raw
Chamfer, while non-rigid ICP reduces it to approximately 18--20\%. Under jitter,
all geometric distances are affected by local noise, but the clean and rigid
perturbation settings show that registration itself can remove substantial
inter-subject variation.

\section{FaceScape Adaptation and Forgetting}
\label{app:faceverse_adaptation}

\subsection{FaceScape Target-Domain Construction}
\label{app:faceverse_construction}

FaceScape~\cite{facescape} is used as an external target domain to assess how the
REMESH-trained metric transfers to a mesh family unseen during source training.
No FaceScape identities are used for REMESH training, checkpoint selection, or
hyperparameter tuning.

We consider two FaceScape evaluation settings. The first is a
\emph{full-neutral} setting, where each subject is represented by one neutral-pose
mesh. We select files matching \texttt{*\_01.npz}, yielding 110 subjects and one
mesh per subject:
\begin{equation}
N_{\mathrm{subjects}} = 110,
\qquad
N_{\mathrm{pairs}} = \binom{110}{2} = 5995.
\label{eq:faceverse_subject_pairs}
\end{equation}
Each mesh has approximately 10k vertices and 20k faces, with DiffusionNet
operators precomputed. Pairwise latent distances are compared against a
FaceScape subject-level distance matrix computed from clean FaceScape detail
meshes using mean per-vertex Euclidean distance.

The second is a harder \emph{cross-topology} setting. Each FaceScape identity is
represented by two mesh realizations:
\begin{itemize}
    \item \textbf{Downsampled FaceScape}: the original FaceScape mesh reduced to
    approximately 10k vertices.
    \item \textbf{Remesh10k}: an independently remeshed version of the same
    subject, also with approximately 10k vertices.
\end{itemize}
Both variants include precomputed DiffusionNet operators. The downsampled
FaceScape meshes contain on average 10,152 vertices and 20,000 faces, while the
remeshed variants contain on average 10,113 vertices and 20,000 faces.

The remeshing pipeline samples points uniformly from the source surface,
estimates and orients normals, performs Poisson surface reconstruction, crops the
result using a scaled source bounding box, and finally applies quadric decimation
to obtain approximately 20,000 faces. This creates a two-realization
cross-topology FaceScape benchmark with 110 subjects, 220 mesh samples, 5,995
subject pairs, and 11,990 mesh pairs.

\subsection{FaceScape Perturbation and Chamfer Protocols}
\label{app:faceverse_perturbation_protocol}

We evaluate FaceScape robustness under five synthetic perturbation scenarios:
clean, jitter, translation, rotation, and mixed. For the fixed-noise experiment,
we use $\sigma=0.05$. Rotation uses a maximum angle of $1+20\sigma$ degrees,
giving $2^\circ$ at $\sigma=0.05$. Translation uses per-axis standard deviation
$0.002+0.05\sigma$, giving $0.0045$ at $\sigma=0.05$. The sigma sweep evaluates
\begin{equation}
\sigma \in \{0.00, 0.02, 0.05, 0.10, 0.15, 0.20\}.
\label{eq:faceverse_sigma_sweep}
\end{equation}

The Chamfer baseline is evaluated both without ICP and, in a separate protocol,
with post-perturbation ICP. The no-ICP protocol uses raw Chamfer distances
computed directly between the evaluated meshes. The ICP variant aligns each
perturbed pair before Chamfer computation using 2048 ICP points, maximum
correspondence distance 0.05, and 20 ICP iterations.

\subsection{Full-Neutral Zero-Shot Evaluation}
\label{app:faceverse_full_neutral}

We first evaluate the frozen REMESH-trained encoder on the full-neutral
FaceScape setting. This setting tests whether the learned metric transfers any
identity-ranking structure to a new mesh family when topology variation is not
introduced.

\begin{table}[h]
\centering
\small
\caption{
FaceScape full-neutral zero-shot evaluation. The REMESH-trained model transfers
nontrivial identity-ranking structure, but Chamfer remains stronger in this
within-topology setting.
}
\label{tab:faceverse_full_neutral_app}
\begin{tabular}{lccc}
\toprule
Scenario & Latent Spearman & Chamfer Spearman & $\Delta$ \\
\midrule
clean & 0.531864 & 0.605202 & -0.073337 \\
mixed & 0.476996 & 0.638449 & -0.161453 \\
\bottomrule
\end{tabular}
\end{table}

The latent metric achieves a clean Spearman correlation of $0.532$, indicating
nontrivial zero-shot transfer rather than out-of-domain collapse. However,
Chamfer remains stronger in this full-neutral within-topology setting, reaching
$0.605$ in the clean case and $0.638$ under mixed perturbations. Thus, this
experiment should be interpreted as evidence of partial transfer, not as a
zero-shot win for the learned metric.

\subsection{Cross-Topology Zero-Shot Evaluation}
\label{app:faceverse_xtopo_zeroshot}

We next evaluate zero-shot transfer in the harder FaceScape cross-topology
setting, where each subject is represented by both a downsampled FaceScape mesh
and an independently remeshed 10k-vertex variant. This introduces target-domain
topology variation that is not present in the full-neutral setting.

\begin{table}[h]
\centering
\small
\caption{
FaceScape cross-topology zero-shot evaluation. This setting exposes a strong
domain shift for the REMESH-trained latent metric.
}
\label{tab:faceverse_xtopo_zeroshot_app}
\begin{tabular}{lccc}
\toprule
Scenario & Latent Spearman & Chamfer Spearman & $\Delta$ \\
\midrule
clean & 0.114807 & 0.598032 & -0.483225 \\
\bottomrule
\end{tabular}
\end{table}

In this harder cross-topology setting, the zero-shot latent metric is
insufficient: Spearman correlation drops to $0.115$, while Chamfer reaches
$0.598$. This result should not be presented as a positive zero-shot outcome.
Instead, it shows that transferring cross-topology robustness from REMESH to a
different mesh family is a substantial domain-shift problem and motivates the
few-shot FaceScape adaptation experiments below.

\subsection{Few-Shot FaceScape Fine-Tuning}
\label{app:faceverse_finetuning}

To assess whether the FaceScape cross-topology gap is due to domain shift, we
fine-tune the REMESH-trained checkpoint on increasing numbers of FaceScape
identities. Fine-tuning is performed in the FaceScape cross-topology setting
using the FaceScape identity distance matrix as supervision. Evaluation uses
held-out subjects and subject-pair mean aggregation.

One representative 75-shot mixed-augmentation run uses 150 symlinked mesh
entries for fine-tuning and evaluates on 35 held-out subjects, 70 samples, and
1,190 subject pairs. Table~\ref{tab:faceverse_75shot_app} reports the strongest
stored held-out result from this branch.

\begin{table}[h]
\centering
\small
\caption{
Representative 75-shot FaceScape held-out cross-topology evaluation. The adapted
model outperforms Chamfer in the clean setting and under moderate translation.
}
\label{tab:faceverse_75shot_app}
\begin{tabular}{lccc}
\toprule
Scenario & Latent Spearman & Chamfer Spearman & $\Delta$ \\
\midrule
clean & 0.537433 & 0.403222 & +0.134211 \\
translation $\sigma=0.10$ & 0.485680 & 0.416691 & +0.068990 \\
\bottomrule
\end{tabular}
\end{table}

After few-shot target-domain adaptation, the learned metric becomes stronger than
Chamfer in the held-out FaceScape cross-topology setting. In the clean case, it
achieves Spearman correlation $0.537$ compared with $0.403$ for Chamfer. Under
moderate translation, it remains stronger, with $0.486$ versus $0.417$.

These results suggest that the weak zero-shot cross-topology performance is
primarily a domain-shift issue rather than a fundamental inability of the learned
metric to represent FaceScape geometry. However, the adapted model is not
uniformly robust: under stronger jitter perturbations, performance degrades and
Chamfer can overtake the learned metric. Thus, FaceScape adaptation improves
target-domain cross-topology ranking, but does not fully solve robustness to
local geometric noise.

\subsection{Source-Domain Forgetting After FaceScape Fine-Tuning}
\label{app:faceverse_forgetting}

Target-domain adaptation improves FaceScape performance, but may reduce the
source-domain invariances learned on REMESH. To assess this, we re-evaluate the
FaceScape-adapted checkpoints on the original REMESH clean cross-topology
protocol. The evaluation uses 100 held-out REMESH identities, 600 meshes, and
148,500 ordered mesh pairs. Since Chamfer does not depend on the learned
checkpoint, we only recompute the latent metric and compare it against the
pre-fine-tuning REMESH model.

\begin{table}[h]
\centering
\small
\caption{
Source-domain performance after FaceScape fine-tuning. The table reports clean
REMESH cross-topology Spearman correlation before and after adapting to
FaceScape. The mean topology-pair score averages Spearman correlations across
ordered topology pairs. Lower post-fine-tuning scores indicate source-domain
forgetting.
}
\label{tab:faceverse_finetune_forgetting_app}
\begin{tabular}{lccc}
\toprule
Model & Overall Spearman & Mean topology-pair Spearman & $\Delta$ mean \\
\midrule
Pre fine-tune REMESH & \textbf{0.751} & \textbf{0.764} & -- \\
FaceScape shot75 clean-biased & 0.595 & 0.637 & -0.128 \\
FaceScape shot75 mixed-augmentation & 0.519 & 0.587 & -0.177 \\
\bottomrule
\end{tabular}
\end{table}

The results show clear source-domain forgetting. The clean-biased FaceScape
checkpoint drops from $0.764$ to $0.637$ mean topology-pair Spearman, while the
mixed-augmentation checkpoint drops further to $0.587$. This indicates that
FaceScape adaptation improves target-domain behavior but does not preserve the
full cross-topology robustness learned on REMESH.

To localize where forgetting occurs, we compute the relative percentage drop for
each ordered topology pair:
\begin{equation}
\mathrm{Drop}_{A \rightarrow B}
=
\frac{
\rho_{\mathrm{pre}}(A \rightarrow B)
-
\rho_{\mathrm{post}}(A \rightarrow B)
}{
\rho_{\mathrm{pre}}(A \rightarrow B)
}
\times 100.
\end{equation}

\begin{table*}[!htbp]
\centering
\scriptsize
\caption{
Relative source-domain forgetting after FaceScape fine-tuning. Each cell reports the percentage drop in REMESH clean cross-topology Spearman correlation relative
to the pre-fine-tuning REMESH checkpoint. Rows indicate source topology and columns indicate target topology. Higher percentages indicate stronger
forgetting.
}
\label{tab:faceverse_finetune_forgetting_matrix_app}
\resizebox{\textwidth}{!}{%
\begin{tabular}{cc}
\begin{tabular}{lrrrrrr}
\toprule
\multicolumn{7}{c}{\textbf{Shot75 clean-biased}} \\
\midrule
Source $\backslash$ Target
& crop & down8k & noisy & original & remesh & up60k \\
\midrule
crop
& 8.2\% & 53.7\% & 44.6\% & 44.4\% & 23.8\% & 43.1\% \\
down8k
& 41.1\% & 3.6\% & 5.7\% & 3.4\% & 14.7\% & 5.5\% \\
noisy
& 35.2\% & 7.6\% & 6.2\% & 6.3\% & 5.6\% & 5.2\% \\
original
& 35.9\% & 5.3\% & 6.6\% & 5.6\% & 3.3\% & 2.7\% \\
remesh
& 18.4\% & 20.8\% & 10.2\% & 7.0\% & 4.8\% & 11.7\% \\
up60k
& 34.2\% & 7.4\% & 6.0\% & 3.0\% & 8.7\% & 4.4\% \\
\bottomrule
\end{tabular}
&
\begin{tabular}{lrrrrrr}
\toprule
\multicolumn{7}{c}{\textbf{Shot75 mixed-augmentation}} \\
\midrule
Source $\backslash$ Target
& crop & down8k & noisy & original & remesh & up60k \\
\midrule
crop
& 11.4\% & 68.2\% & 58.8\% & 59.3\% & 39.7\% & 58.5\% \\
down8k
& 55.9\% & 5.1\% & 8.4\% & 5.2\% & 18.4\% & 7.8\% \\
noisy
& 48.5\% & 12.1\% & 8.9\% & 9.2\% & 9.2\% & 9.0\% \\
original
& 49.8\% & 7.8\% & 8.7\% & 7.1\% & 6.2\% & 5.0\% \\
remesh
& 33.5\% & 24.8\% & 12.8\% & 9.7\% & 6.5\% & 14.3\% \\
up60k
& 49.2\% & 9.9\% & 8.7\% & 5.0\% & 11.4\% & 6.1\% \\
\bottomrule
\end{tabular}
\end{tabular}
}
\end{table*}

The forgetting pattern is not uniform. The largest degradation occurs for
topology pairs involving cropped meshes, where the drop often exceeds $40\%$ and
can reach more than $60\%$ for the mixed-augmentation checkpoint. In contrast,
pairs among non-cropped variants such as original, noisy, downsampled, and
upsampled meshes degrade much less, often in the range of $3$--$12\%$ for the
clean-biased checkpoint.

This suggests that FaceScape fine-tuning shifts the representation toward the
target-domain topology statistics and weakens the invariances learned for the
most severe REMESH topology changes, especially partial-support cases. Therefore,
the FaceScape-adapted checkpoint should not be interpreted as a new universal
model. These results motivate future adaptation strategies that explicitly
control forgetting, such as mixed REMESH--FaceScape rehearsal, distillation from
the REMESH checkpoint, partial freezing of the encoder, or multi-domain training.

\FloatBarrier
\newpage
\begin{enumerate}
    \item For all authors: this draft is preliminary and the checklist must be completed before submission.
    \item The main claims are stated in the introduction, results, and conclusion.
    \item The experimental protocol and current aggregate results are described in the main text.
    \item Code, data paths, hyperparameters, and exact training details should be finalized before submission.
    \item Limitations are discussed in the main text.
    \item Broader impacts should be added if required by the final NeurIPS 2026 instructions.
\end{enumerate}


\end{document}